% The odd case, and the sweep.
%
% This was Appendix A of the paper.  It records the case analysis that the
% trichotomy at the end of Step 4 replaced -- which lambda actually occur at odd
% a, rather than merely that each one is harmless -- and the computational sweep
% that checks both.  Nothing in the paper depends on it, which is why it was
% moved here.  Kept because it is how the step was first proved, and because
% Remark 27 in it records a claim of ours that turned out to be false.
%
% It is not a standalone document; the macros are the paper's.

\appendix

\section{The odd case, and the sweep}\label{app:odd}

The trichotomy at the end of Step 4 settles the odd case together with the even
one, for every $\lambda$ at once. This appendix records the finer analysis it
replaced---which says not merely that each $\lambda \ne -1$ is harmless but which
$\lambda$ occur---and the computational sweep that checks both against every
rank-$8$ spindle to $a = 80$. Nothing in the paper depends on either.

\subsection*{Which $\lambda$ occur at odd $a$}
For odd $a$ Lemma \ref{lem:E} says nothing, and two sub-cases exhaust what occurs.

\emph{(a) $\lambda$ rational.} Write $\lambda = \alpha/\beta$ with
$\alpha, \beta$ coprime integers, $\beta \ge 1$. A rational $\lambda$ makes
$m = \lvert v \rvert^2$ rational, so $m$ is a rational squared norm of a non-zero
point of $\Lat$, hence $m \ge 1/9$ by Section \ref{sec:arith}. Rearranged,
\eqref{eq:star} reads $m\big[a(\alpha+\beta)^2 - \alpha\beta\big] = a$. Both
signs end up bounding the same quantity, $\lvert\alpha\rvert - \lvert\beta\rvert$,
but for different reasons.

\emph{Negative $\lambda$}, $\lambda = -p/q$ with $p,q \ge 1$, $d = q-p$. Here
\eqref{eq:star} reads $a d^2 + pq = a/m$. With $pq \ge 1$,
\[
  d^2 < 1/m \le 9, \qquad\text{so}\qquad \lvert d \rvert \le 2 ,
\]
and $d = 0$ is $\lambda = -1$; $\lvert d \rvert = 3$ is impossible, not merely
unobserved. If $\lvert d \rvert = 1$ then exactly one of $p,q$ is even, so $2$ and
hence both primes divide one of $r,s$, while the other is an odd multiple of $v$
and Corollary \ref{cor:even} keeps its colour non-zero: $\varphi(u) = 0 + c$ with
$c \ne 0$ once more. (Here $a$ is odd, so it is the first half of the corollary
that applies, not the second.) If $\lvert d \rvert = 2$ then $p$ and $q$ are coprime and differ by $2$, so
both are odd; $4a + pq = a/m$ forces $m < 1/4$, and $1/9$ is the only squared
norm below $1/4$ (Section \ref{sec:arith}), so $m = 1/9$ exactly and
$pq = 5a$. Then $r = qv$ and $s = -pv$ with $\lvert v \rvert^2 = 1/9$ and $p,q$
odd, so $\pi_i$ divides neither, multiplication by them is invertible mod $\pi_i$,
and $\gamma_i(r) = 0 \iff \pi_i \mid v \iff \gamma_i(s) = 0$.

\begin{lemma}\label{lem:F}
In that case $\pi_i$ does not divide $v$.
\end{lemma}

\begin{proof}
If $v = \pi w$ with $w \in \Lat$ then, $\pi$ being real,
$\lvert v \rvert^2 = \pi^2 \lvert w \rvert^2$, so
$\lvert w \rvert^2 = 1/(9\pi^2) = (58 \mp 10\sqrt{33})/144$ and \eqref{eq:norm}
demands $A^2 + 33B^2 + 3C^2 + 11D^2 = 58$. The congruences force
$A \equiv B \equiv C \equiv D \pmod 2$: all even makes the form $\equiv 0 \pmod 4$
against $58 \equiv 2$; all odd makes every square $\equiv 1 \pmod 8$, so the form
is $\equiv 1 + 33 + 3 + 11 = 48 \equiv 0 \pmod 8$, again against $58 \equiv 2$.
Neither is possible.
\end{proof}

So both colours are equal and non-zero, and the vector imposes the same diagonal
condition $(c,c)$ that the shell displacements of Step 2 impose. Nothing new.
(Two shorter routes reach the same place. Corollary \ref{cor:even} gives
$c \ne 0$ at once, since $c = 0$ would put $v$ in $\pi_i\Lat$; and $c = 0$ would put
$v$ in $2\Lat$, forcing $\lvert v/2 \rvert^2 = 1/36$, which Section
\ref{sec:arith} records as unattained---congruence (ii) refuses its only candidate,
$(\pm2,0,0,0)$. That second route is the one the accompanying code takes.
Lemma \ref{lem:F} needs neither, and proves the stronger statement that neither
prime divides $v$.)

\emph{Positive $\lambda$}, $\alpha,\beta \ge 1$. This sign occurs and is not
covered by the bound above: $\alpha\beta \ge 1$ now, so the rearrangement gives
$a\big[m(\alpha+\beta)^2 - 1\big] = m\alpha\beta$, and $4a > 1$---which
$\sqrt{4a-1}$ needs---is exactly $m(\alpha-\beta)^2 < 1$, so $(\alpha-\beta)^2 < 9$
again. The dichotomy is parity, not size. If $\alpha,\beta$ have opposite parity,
one is even, both primes divide the corresponding one of $r,s$, and the other
colour is again non-zero by Corollary \ref{cor:even}. If both are odd then
$\gamma_i(r) = \gamma_i(s) = \gamma_i(v)$ and the pair is the diagonal $(c,c)$ of
Step 2, with $c \ne 0$ by Corollary \ref{cor:even}. Nothing new either way. The
case is not vacuous: $a = 7/19$ is odd and of rank $8$, and
$v = \tfrac56 + \tfrac{i\sqrt3}{6}$, of squared norm $7/9$, gives a cross unit
vector with $\lambda = 1$.

\emph{(b) $\lambda$ irrational.} Then $\lvert v \rvert^2$ is irrational and the
bound of (a) is unavailable; the three cases above need none of it. For the record,
irrational $\lambda$ occurs only at even $a$ in every case computed
($a = 4, 10, 16, 28$, over all rank-$8$ spindles to $a = 80$), but nothing rests on
that now: it was the last computed-only step in the proof of
Theorem~\ref{thm:law}, and the trichotomy removes it.

\subsection*{What the sweep checks}
The accompanying code checks Step 4 and the analysis above against the actual
cross unit vectors of every rank-$8$ spindle to $a = 80$, with zero violations,
and both halves of
Corollary \ref{cor:even} on all $552$ cross vectors with $\lambda \ne -1$ to
$a = 60$. Each prime divides $r$ or $s$ in every one of them---not only at even
$a$, where the corollary asks for it, but at odd $a$ as well---with exactly two
exceptions, $a = 39$ and $a = 51$. Both are the $\lvert d \rvert = 2$ case above,
where $p$ and $q$ are odd and neither $r$ nor $s$ is divisible; there Lemma \ref{lem:F}, or the trichotomy, applies instead. The trichotomy is checked
on the same $552$ directly: where $4 \nmid a$, no prime divides \emph{both} $r$
and $s$, again with zero counterexamples.

\begin{remark}
An earlier form of this analysis asserted that every rational $\lambda$ is a ratio of
\emph{consecutive} integers, $\lvert p - q \rvert = 1$, measured over every
rank-$8$ spindle with $a \le 35$. That is false: at $a = 39$ the cross vectors have
$\lambda = -13/15$ and $-15/13$. The measurement had stopped four short of its own
counterexample. The bound $\lvert d \rvert \le 2$ above is a proof, and
$\lvert d \rvert = 2$ occurs exactly when $pq = 5a$ with $m = 1/9$---which to
$a = 80$ means $a = 39$ and $a = 51$, the two the sweep singles out, since
$a = p(p+2)/5$ gives rank $4$ at $p = 3$ and $5$ and lands past $80$ at $p = 23$.
The positive-$\lambda$ branch refutes the old assertion a second time, from the other
side: $\lambda = 1$ has $\lvert \alpha - \beta \rvert = 0$.
\end{remark}
